% !TeX program = lualatex
% Compile twice: lualatex open_problems_500.tex
% All references and prose are embedded; no BibTeX database or external inputs.
\documentclass[11pt,a4paper]{article}
\usepackage[margin=24mm,headheight=14pt]{geometry}
\usepackage{fontspec}
\setmainfont{Latin Modern Roman}
\setsansfont{Latin Modern Sans}
\setmonofont{Latin Modern Mono}
\usepackage{amsmath,amssymb,mathtools}
\usepackage{unicode-math}
\setmathfont{Latin Modern Math}
\usepackage{microtype}
\usepackage{enumitem,xurl,needspace}
\usepackage[dvipsnames]{xcolor}
\usepackage{hyperref}
\hypersetup{unicode=true,colorlinks=true,linkcolor=MidnightBlue,urlcolor=MidnightBlue,
 pdftitle={500 Mathematical Problem Statements},pdfauthor={Source-annotated compilation},bookmarksnumbered=true}
\usepackage{bookmark}
\usepackage{tocloft}
\setlength{\cftsecnumwidth}{3em}
\cftsetpnumwidth{2.5em}
\cftsetrmarg{4em}
\usepackage{titlesec}
\titleformat{\section}{\Large\bfseries\raggedright}{\thesection}{0.7em}{}
\usepackage{fancyhdr}
\pagestyle{fancy}\fancyhf{}
\fancyhead[L]{\small\sffamily Mathematical problem statements}
\fancyhead[R]{\small Source edition 22}
\fancyfoot[C]{\thepage}
\setlength{\parindent}{0pt}
\setlength{\parskip}{5pt plus 1pt}
\setlength{\emergencystretch}{3em}
\setcounter{tocdepth}{1}
\setcounter{secnumdepth}{1}
\setlist[itemize]{leftmargin=3em,itemsep=5pt,topsep=4pt,parsep=0pt}
\allowdisplaybreaks
\newcommand{\N}{\mathbb{N}}
\newcommand{\Z}{\mathbb{Z}}
\newcommand{\Q}{\mathbb{Q}}
\newcommand{\R}{\mathbb{R}}
\newcommand{\C}{\mathbb{C}}
\newcommand{\F}{\mathbb{F}}
\newcommand{\A}{\mathbb{A}}
\newcommand{\PP}{\mathbb P}
\newcommand{\E}{\mathbb{E}}
\newcommand{\eps}{\varepsilon}
\newcommand{\dd}{\,\mathrm d}
\newcommand{\sref}[2]{\hyperlink{source:#1:#2}{[S#2]}}
\newcommand{\eref}[2]{\hyperlink{extra:#1:#2}{[E#2]}}
% Turn the next switch off for a shorter reading copy. The quotations preserve
% every exactTarget field, including qualifications omitted from a short summary.
\newif\ifshowcatalogue
\showcataloguetrue
\newcommand{\cataloguescope}[1]{\ifshowcatalogue
 \paragraph{Exact scope in the supplied catalogue.}
 \begin{quote}\small #1\end{quote}\fi}

% Standalone research notebook, original catalogue problem 192.
% Compile twice with: lualatex 192.tex
\numberwithin{equation}{section}
\hypersetup{pdftitle={Research attempt -- problem 192},pdfauthor={Research notebook; original compilation retained}}
\fancyhead[L]{\small\sffamily Research notebook}
\fancyhead[R]{\small\sffamily Problem 192 | 27 September 2026}
\begin{document}
\setcounter{section}{191}
% ================================================================
% problemId: problem.grothendieck-serre-conjecture-on-principal-bundles-over-regular-local-rings
% source releaseRank: 192; source releaseStatus: open_with_solved_subcases
% exactTarget SHA-256: e92dbd17cd4f7f9995bc1d79926910a07c239174233c4079ddcea2e876c992cb
\clearpage
\section{\texorpdfstring{Grothendieck–Serre Conjecture on Principal Bundles over Regular Local Rings}{Grothendieck–Serre Conjecture on Principal Bundles over Regular Local Rings}}\label{192}
\subsection{Definitions and mathematical statement}
A Noetherian local ring $(R,\mathfrak m)$ is regular when
$\dim R=\dim_{R/\mathfrak m}\mathfrak m/\mathfrak m^2$; it is a domain here, with fraction field $K$. A reductive $R$-group scheme is smooth affine with connected reductive geometric fibers. A $G$-torsor is a principal homogeneous space locally trivial for the \"etale topology. The assertion is
\[
 \ker\bigl(H^1_{\mathrm{et}}(R,G)\to H^1_{\mathrm{et}}(K,G)\bigr)=\{1\},
\]
where these are pointed sets and the kernel is the preimage of the trivial torsor. Thus generic triviality forces triviality over $R$. The exact statement does not impose a field inside $R$, excellence, isotropy, or a constant group scheme.
\par\smallskip\textit{Formulation and concept references:} \sref{192}{1}.
\cataloguescope{For every commutative Noetherian regular local ring R, with fraction field K=Frac(R), and every reductive R-group scheme G (smooth affine with connected reductive geometric fibers), every G-torsor over Spec(R) that becomes trivial over Spec(K) is already trivial over Spec(R). Equivalently, the natural map of nonabelian pointed sets H\textasciicircum{}1\_et(R,G) -> H\textasciicircum{}1\_et(K,G) has trivial kernel. Regular local rings here are domains. No field-containing, excellence, finite-type, unramified, quasi-split, isotropic or constant-group restriction is imposed.}
\subsection{Short English statement}
If a principal bundle under a reductive group over a regular local ring is trivial over the fraction field, must it already be trivial over the ring?
% BEGIN RESEARCH ADDITION ATTEMPT-192
\subsection{Research attempt}

\paragraph{Current frontier.}
The cited survey distinguishes the full regular-local conjecture from field-containing and unramified cases \sref{192}{1}. A recent theorem covers unramified regular semilocal rings with totally isotropic reductive groups, still not the unrestricted pair $(R,G)$ here \eref{192}{1}. The discrete-valuation-ring case is an established input, as recorded in the survey. We investigate whether codimension-one trivializations can be synchronized into one generic section. This separates the easy normality extension from the genuinely nonabelian patching problem.

\paragraph{Attack route.}
There are three plausible directions: descend a trivialization from a completion, use approximation in $G(K)$ to correct generic denominators, or apply a geometric presentation/patching theorem. Completion descent requires compatible descent data, not just a point after base change. We pursue the second route and identify an exact simultaneous-coset condition; the last route would have to prove that condition for the torsor-derived data in the remaining generality.

\paragraph{Attempt.}
Let $P$ be a generically trivial right $G$-torsor. Since $G$ is affine, affineness descends along an \"etale trivializing cover, so $P=\operatorname{Spec}B$ is affine over $R$. Choose $s\in P(K)$.

\textbf{Hartogs extension lemma.} If this \emph{same} point $s$ extends to $P(R_{\mathfrak p})$ for every height-one prime $\mathfrak p$, then it extends uniquely to $P(R)$.
A regular local domain is normal, and a Noetherian normal domain satisfies
\begin{equation}\label{eq192:intersection}
 R=\bigcap_{\operatorname{ht}\mathfrak p=1}R_{\mathfrak p}\subset K.
\end{equation}
The coordinate homomorphism $B\to K$ corresponding to $s$ therefore has image in $R$, proving existence; injectivity $R\hookrightarrow K$ proves uniqueness. In dimension zero the ring is already a field and the assertion is immediate. Equation~\eqref{eq192:intersection} is the standard height-one description of a normal Noetherian domain \eref{192}{2}.

Because $B$ is finitely presented, $s$ extends over $R[1/f]$ for some nonzero $f$. Only finitely many height-one primes contain $f$. For each of those primes, the discrete valuation ring theorem supplies a local point $t_{\mathfrak p}\in P(R_{\mathfrak p})$. Since $P_K$ is a torsor, there is a unique $g_{\mathfrak p}\in G(K)$ such that
\[
 t_{\mathfrak p}=s\cdot g_{\mathfrak p}.
\]
For primes away from the bad set put $g_{\mathfrak p}=1$. The set of corrections making $s\cdot g$ integral at $\mathfrak p$ is exactly
\[
 g_{\mathfrak p}G(R_{\mathfrak p})\subset G(K).
\]
Indeed, an integral point trivializes the torsor over $R_{\mathfrak p}$, so every other integral point differs by an element of $G(R_{\mathfrak p})$.

Thus a sufficient and necessary patching condition for this torsor and this generic section is
\begin{equation}\label{eq192:patch}
 \bigcap_{\operatorname{ht}\mathfrak p=1}
          g_{\mathfrak p}G(R_{\mathfrak p})\ne\varnothing.
\end{equation}
For $g$ in the intersection, $s\cdot g$ extends everywhere in codimension one and the Hartogs lemma gives a section of $P$ over $R$, hence triviality. Conversely any global section determines such a $g$. Notice the order $g_{\mathfrak p}^{-1}g\in G(R_{\mathfrak p})$; swapping factors would change a nonabelian condition.

This is not a claim that arbitrary collections of local cosets can always be synchronized. The collection here is constrained by coming from a single torsor over $R$. Proving \eqref{eq192:patch} just for those collections is enough; making it an unrestricted approximation assertion would be stronger and could be false.

\textbf{Abelian calibration.} For $G=\mathbf G_m$, regular local factoriality permits an explicit correction. Write the prescribed valuations at the finitely many bad primes as $e_{\mathfrak p}=v_{\mathfrak p}(g_{\mathfrak p})$, choose generators $\pi_{\mathfrak p}$ of these height-one primes, and take
\[
 g=\prod_{\mathfrak p\text{ bad}}\pi_{\mathfrak p}^{e_{\mathfrak p}}\in K^\times.
\]
It has the required valuation at every height-one prime, including zero at all others; hence it lies in every coset. The same argument works coordinatewise for a split torus. This rederives an elementary special case: invertible modules over a local ring are free. For $\mathrm{GL}_r$, torsors are finite locally free rank-$r$ modules and local freeness likewise makes every torsor trivial. These calibrations show that nonabelian structural constraints, not the normality intersection itself, contain the general difficulty.

\paragraph{Stress test.}
The false quantifier swap would be
\[
 \bigl[\forall\mathfrak p\ \exists t_{\mathfrak p}\in P(R_{\mathfrak p})\bigr]
 \quad\Longrightarrow\quad
 \bigl[\exists s\in P(K)\ \forall\mathfrak p\ (s\text{ integral at }\mathfrak p)\bigr].
\]
The Hartogs lemma proves only the conclusion after the single section has been found. Weak approximation at finitely many primes may introduce new bad denominators elsewhere, whereas \eqref{eq192:patch} includes integrality at \emph{all} the other primes. A proof that ignores them is incomplete. Nor does an $R$-torsor point after completion descend automatically to $R$ without a descent argument.

No excellence, coefficient field, isotropy or unramifiedness was inserted into the Hartogs calculation. Such hypotheses do appear in available theorems that might supply the missing patching; they cannot be dropped when applying those theorems. The current attempt does not produce a patching theorem in the arbitrary ramified and anisotropic setting.

\paragraph{Outcome.}
\textbf{Outcome: Conditional reduction.}
The full target for a given torsor is reduced to a precise intersection of torsor-derived cosets in $G(K)$. The extension step and the split-torus calibration are proved, and the difference between local existence and a simultaneous generic correction is explicit. The unresolved ingredient is nonabelian synchronization in the full regular-local generality, not a further application of normality.

% END RESEARCH ADDITION ATTEMPT-192
\subsection{Sources}
\begin{itemize}\raggedright
\item[\textnormal{[S1]}]\hypertarget{source:192:1}{} Problems about torsors over regular rings, with an appendix by Yifei Zhao; arXiv2201.06424v4,5May2025; HTML internal29March2023.
\url{https://arxiv.org/pdf/2201.06424v4}.
{\small\textit{Catalogue formulation source.}}
% BEGIN RESEARCH ADDITION REFERENCES-192
\item[\textnormal{[E1]}]\hypertarget{extra:192:1}{} K\k{e}stutis \v{C}esnavi\v{c}ius, \emph{Unramified Grothendieck--Serre for isotropic groups}, arXiv:2311.08660, main theorem and its unramified/totally isotropic hypotheses.
\url{https://arxiv.org/abs/2311.08660}.
{\small\textit{Research reference; consulted 27 September 2026.}}
\item[\textnormal{[E2]}]\hypertarget{extra:192:2}{} The Stacks Project Authors, \emph{The Stacks Project}, Algebra, Lemma on a Noetherian normal domain as the intersection of its height-one localizations (Tag 031T), and regular local factoriality.
\url{https://stacks.math.columbia.edu/tag/031T}.
{\small\textit{Research reference; consulted 27 September 2026.}}
% END RESEARCH ADDITION REFERENCES-192
\end{itemize}

\end{document}
